\documentclass{article}

\usepackage[dblblindworkshop]{neurips_2026}

\workshoptitle{I Can't Believe It's Not Better (ICBINB): Failure Modes of AI in Biology}

\usepackage[utf8]{inputenc}
\usepackage[T1]{fontenc}
\usepackage[hidelinks]{hyperref}
\usepackage{url}
\usepackage{booktabs}
\usepackage{array}
\usepackage{amsfonts}
\usepackage{amsmath}
\usepackage{nicefrac}
\usepackage{microtype}
\usepackage{xcolor}
\usepackage{graphicx}
\usepackage{placeins}

\title{Compositional Separation Predicts What Protein Language Model
Steering Can and Cannot Do}

\author{Anonymous Authors}

\begin{document}

\maketitle

\begin{abstract}
Activation steering edits a protein language model's residual stream at
inference time and can change its generated sequences without retraining, but a
higher proxy score does not by itself establish control of a biological
property, and the field lacks an \emph{a priori} rule for which properties a
given steering method can move at all. We attempted to steer nine sequence
properties in ESM2-650M under one fixed harness---mean-pooled per-layer
difference-of-means vectors added uniformly at every position, evaluated by
single-shot masked filling---and one pre-registered six-criterion acceptance
protocol. A single quantity measured before any generation, the mean amino-acid
\emph{compositional separation} between the low- and high-label groups used to
build the steering vector, ordered every outcome: it predicts both whether a
target is steerable and how large the achievable effect is, whereas the
proxy's correlation with the true label does not. A binding proxy validated at
held-out $r=0.80$ steered nothing (separation $0.003$), while a catalytic proxy
at $r=0.22$ (separation $0.02$--$0.05$) passed on all three seeds; effect size
grew with separation from a flat null to $72\%$ of the natural property gap.
Separation is necessary but not sufficient: two well-separated targets still
failed, one because the model does not act on the direction and one because the
property is \emph{positional}. That last case exposes a structural limit:
uniform-addition steering can move aggregate compositional properties but
cannot install positionally-localized features, regardless of proxy quality.
All endpoints are computational; no wet-lab assay was run.
\end{abstract}

\section{Problem}
\label{sec:problem}

Protein language models learn statistical structure from large sequence
collections and support prediction and generation \citep{lin2023esm}.
Activation steering modifies their internal representations at inference time:
a direction is estimated from the difference between activations for high-label
and low-label examples, then added to the residual stream during generation.
Prior work applied this method to protein generation and optimization
\citep{huang2025steering}.

Two problems sit between a steering run and a useful conclusion. First, a
generated sequence can receive a higher proxy score for reasons unrelated to
the claimed property---by repeatedly inserting residues that the scoring rule
rewards, or after low-complexity outputs have been silently dropped from the
comparison. This half of the problem is an evaluation-validity question, and
prior negative-results work has catalogued how each such check can overturn an
apparent success. Second, and less studied, is a \emph{selection} question:
given a fixed steering method, which target properties can it move at all, and
can that be predicted before spending any compute? A practitioner who can only
learn steerability by running the full harness on each candidate has no way to
budget effort or to know when a null is a property of the method rather than of
the target.

This paper reports what a single arc of nine steering attempts on ESM2-650M
revealed about the second question. The central finding is a predictor and a
structural limit, and both are, in the venue's sense, failure results: most
well-separated candidates still fail, and the method is provably unable to
reach an entire class of biologically important properties.

\section{Proposed Approach}
\label{sec:approach}

\subsection{One fixed harness}

Every target uses the same intervention and evaluation, so that outcomes are
comparable across properties. The model is ESM2-650M. Generation uses
single-shot masked filling (\texttt{mask\_fill\_generate}, mask fraction
$0.3$, argmax decode). The steering direction is the per-layer difference
between mean activations of 150 high-label and 150 low-label proteins, added to
all 33 layers uniformly across sequence positions. The control is a
matched-norm random direction. Each run evaluates 150 held-out proteins; a
target-specific compositional proxy scores each generation, and the reported
effect is the paired-bootstrap contrast (10,000 resamples) between the real and
random directions.

The masked-fill evaluation collapses into degenerate low-complexity output at
steering strength $\alpha\!\ge\!0.35$; dose-response is therefore measured on
the non-degenerate grid $\{0.10, 0.15, 0.20, 0.25\}$. This caps the achievable
effect but keeps every comparison inside a valid operating window.

\subsection{A pre-registered acceptance protocol}

To separate real capability from a flattering scalar, each target was judged
against six criteria fixed in advance. A capability gain counts only if it
\textbf{(1)} beats the random-direction control with a real paired-bootstrap
confidence interval; \textbf{(2)} shows a coherent dose-response over at least
three strengths; \textbf{(3)} survives \emph{residue exclusion}---the effect
must remain significant after the two dominant substituted residues are removed
from the score; \textbf{(4)} uses a proxy validated against real experimental
labels \emph{before} the run, asserted in code at $|r|\!\ge\!0.15$;
\textbf{(5)} beats any prior technique for the property; and \textbf{(6)} is
adequately powered ($n\!\ge\!150$). Failing (1), (2), (3), or (4) outright is a
\textsc{kill}; passing most but failing a dose check or (3) is
\textsc{ambiguous}. Criterion (3) is deliberately strict because every proxy
here is a function of amino-acid composition, so steering can raise it simply by
enriching favored residues.

\subsection{A compositional-separation gate, promoted to a predictor}

Before any generation, we measure the mean per-residue amino-acid-composition
$L_2$ distance between the low- and high-label groups from which the vector is
built. This gate was introduced to explain one early null and became the arc's
most predictive number. The claim tested across the nine targets is that this
pre-run separation, not the proxy's label correlation, orders both
steerability and effect size.

\subsection{Targets and data}

Nine properties were attempted, chosen to vary compositional separation and to
contrast intrinsic against relational properties. Transmembrane fraction,
DNA-binding propensity, catalytic turnover, glycosylation density, disulfide
content, signal peptide, zinc-finger and calcium-binding domains were drawn
from reviewed UniProt \citep{uniprot2023}, cross-protein $k_{\mathrm{cat}}$ from
DLKcat \citep{li2022dlkcat}, single-backbone binding from ProteinGym
\citep{notin2023proteingym}, and disorder from DisProt
\citep{quaglia2022disprot}. Proxies use published biophysical scales: mean
Kyte--Doolittle hydropathy for membrane character \citep{kyte1982hydropathy},
Lys$+$Arg fraction for nucleic-acid interfaces \citep{luscombe2001dna}, a
glycine--arginine contrast for the activity--stability tradeoff in turnover
\citep{siddiqui2006cold}, and the TOP-IDP scale for disorder
\citep{campen2008topidp}. An earlier immune-endpoint target using IEDB
\citep{vita2019iedb} and a solubility target using eSol \citep{niwa2009esol}
were stopped by the validity checks and are discussed only where they inform
the selection question.

\section{Observed Outcome}
\label{sec:outcome}

\subsection{Separation orders steerability and effect size}

Table~\ref{tab:scoreboard} lists all nine targets sorted by verdict. Effect
size tracks pre-run compositional separation monotonically, while the proxy's
label correlation does not. The two anchoring rows are the sharpest: the
single-backbone binding target has the strongest-validated proxy in the entire
arc (held-out $r=0.80$) and produces a clean null at every strength, because its
low- and high-label groups are one- and two-residue point mutants of one
backbone and are therefore nearly compositionally identical (separation
$0.003$). Catalytic turnover has the weakest runnable proxy ($r=0.22$) yet
passes on all three seeds. Proxy-validation strength and steering effect are not
even weakly monotonically related.

\FloatBarrier
\begin{table}[htbp]
\centering
\caption{Nine steering targets on ESM2-650M under one fixed harness, sorted by
verdict. Separation is the pre-run compositional $L_2$ distance between vector
groups; proxy $r$ is the held-out label correlation; effect is expressed, where
a natural low$\rightarrow$high proxy gap is defined, as the fraction of that gap
reached at $\alpha=0.25$. Effect size follows separation, not proxy $r$.}
\label{tab:scoreboard}
\small
\begin{tabular}{@{}>{\raggedright\arraybackslash}p{3.5cm}
                    r r
                    >{\raggedright\arraybackslash}p{2.1cm}
                    >{\raggedright\arraybackslash}p{3.4cm}@{}}
\toprule
Target & Sep. & Proxy $r$ & Verdict & Effect / failure mode \\
\midrule
Binding, single-backbone   & 0.003 & 0.80 & \textsc{kill} & flat null (no separation) \\
Zinc-finger domain         & 0.028 & 0.29 & \textsc{kill} & direction anti-correlated \\
Calcium-binding (EF-hand)  & 0.038 & 0.16 & \textsc{kill} & model inert to direction \\
Signal peptide             & 0.035 & 0.74 & \textsc{kill} & positional, not compositional \\
\addlinespace
Disulfide (Cys)            & 0.026 & 0.17 & pass (2/3) & tiny ($\sim\!0.002$) \\
Glycosylation density      & 0.039 & 0.41 & pass (3/3) & tiny ($\sim\!0.003$) \\
DNA-binding propensity     & 0.039 & 0.25 & pass (2/3) & 25\% of natural gap \\
Catalytic turnover ($k_{\mathrm{cat}}$) & 0.02--0.05 & 0.22 & pass (3/3) & moderate \\
Transmembrane fraction     & 0.086 & 0.79 & pass (3/3) & 72\% of natural gap \\
\bottomrule
\end{tabular}
\end{table}
\FloatBarrier

The calibration is visible along the passing rows: binding at separation
$0.003$ gives a null; DNA-binding at $0.039$ moves generations $\sim\!25\%$ of
the way to the binder distribution; transmembrane at $0.086$ moves them
$\sim\!72\%$ of the way to the membrane-protein distribution
($\alpha=0.25$ proxy $-0.363\!\rightarrow\!+0.266$ against a natural gap of
$0.871$). Moving the binding target from single-backbone to cross-protein data
(DNA-binding) lifts separation $\sim\!12\times$ and converts the null into a
significant, residue-robust, dose-responsive effect. The earlier
single-backbone null was thus an artifact of the data geometry, not a property
of binding, and the recovery was predicted from separation before any GPU run.

\subsection{Separation is necessary but not sufficient}

Screening by separation removes the null mode but does not guarantee success.
Two targets with separation as high as the passing band still failed, by
distinct mechanisms. Calcium-binding (separation $0.038$) shows no significant
steering at any safe strength: a compositionally distinct group whose direction
the residual stream simply does not act on. Zinc-finger's cysteine--histidine
direction becomes \emph{anti-correlated} with its proxy after steering. High
separation is a prerequisite, but the model must also respond to the direction,
and the direction must encode the property rather than a confound.

\subsection{The harness steers composition, not position}

The signal-peptide target is the sharpest failure. Its proxy---mean hydropathy
over the N-terminal 30-residue window---is the best-validated of any candidate
in the arc (held-out $r=0.742$), and it produces a textbook dose-response
(N-terminal hydropathy $0.059\!\rightarrow\!0.100\!\rightarrow\!0.152
\!\rightarrow\!0.251$ across the safe grid). It nonetheless
\textsc{kill}s on criterion (3): excluding the two dominant substituted
residues collapses the effect to non-significance. The steering vector raises
N-terminal-window hydropathy purely by enriching bulk hydrophobic residues
across the \emph{whole} sequence; it cannot \emph{concentrate} them at the
N-terminus, which is what defines a signal peptide. Glycosylation density, by
contrast, passes precisely because its proxy is a sequon count summed over the
whole sequence---a genuinely aggregate quantity a mean-pooled vector can nudge.

\section{Reason for Failure}
\label{sec:failure}

\subsection{Why proxy correlation misleads and separation predicts}

A difference-of-means direction can encode only what differs, on average,
between the two groups it is built from. When the groups are near-identical in
composition---single-backbone point mutants---the direction carries almost no
signal, no matter how faithfully the \emph{scoring} function tracks the real
label. Proxy correlation certifies the ruler; separation certifies that there
is something for the vector to push. This is why the best-validated proxy in the
arc produced the cleanest null, and why a weak proxy over compositionally
separated groups produced a replicated pass. Effect size then scales with
separation because a larger mean gap gives the added direction more magnitude to
move generations along.

\subsection{Why well-separated targets still fail}

Separation is a property of the data; steerability additionally requires that
the model's residual stream respond to the direction and that the direction
capture the property rather than a correlate. Calcium-binding fails the first:
the acidic-residue direction is well-separated but inert under this
intervention. Zinc-finger fails the second: its direction moves against its own
proxy. Neither failure is visible from separation or proxy correlation alone,
so the pre-registered head-to-head control (criterion 1) and residue-exclusion
(criterion 3) remain necessary even after the separation gate is cleared.

\subsection{Why an entire property class is unreachable}

The signal-peptide result is not a tuning failure but a structural one. A vector
that is mean-pooled over positions and added identically to every position can
change the \emph{amount} of a residue type in a sequence but not \emph{where}
it goes. Aggregate compositional properties---transmembrane fraction, catalytic
composition, DNA-binding charge, glycosylation density---are functions of that
amount and are reachable. Positionally-localized features---a signal peptide,
and by the same argument active-site geometry, domain boundaries, and
termini-specific motifs---are functions of arrangement and are not. Reaching
them would require position-aware injection, a different intervention, not a
different target. This bounds the class of biological properties this widely
used steering recipe can address.

\subsection{The validity checks that make the predictor trustworthy}

The predictor is only credible because the acceptance protocol rejected
flattering numbers elsewhere in the arc, and each rejection changed a
conclusion. An immune-endpoint proxy correlated $r=0.427$ with an
allele-aggregated binding label but only $0.100$ with measured T-cell response,
and its full-length correlation fell from $+0.379$ under random folds to
$-0.323$ under organism-grouped folds---so it was stopped before generation for
measuring the wrong endpoint. A layer-subset thermostability comparison looked
competitive only because the all-layer arm's outputs had collapsed under the
low-complexity filter, leaving too few evaluable pairs; inside the shared valid
range the subset retained just $43$--$59\%$ of the full effect. A solubility
contrast of $+0.0125$ fell to $+0.00035$, with a confidence interval crossing
zero, once its two dominant substituted residues were excluded. A disorder
effect repeated across three stored runs but survived residue exclusion in only
two. Each check protects a different inference: endpoint validity, evaluable
denominator, composition dependence, and whole-run stability. Only after a
target clears all of them does its separation-versus-effect placement in
Table~\ref{tab:scoreboard} carry weight.

\section{Discussion and Limitations}
\label{sec:limitations}

The useful consequence is that, for this harness, steerability and its ceiling
are predictable before any compute from the compositional separation of the
vector-building groups, and an entire property class (positional features) is
ruled out on structural grounds. The productive next steps are therefore not to
try more properties but to change the intervention---position-aware
injection---and to validate the strongest compositional gain against a
biological assay.

The evidence has firm limits. Every endpoint is computational: no wet-lab or
structural assay establishes that any generated sequence has improved
thermostability, turnover, transmembrane character, binding, or glycosylation;
each ``effect'' is a change in a compositional proxy. All results use one model
(ESM2-650M), one decoder (single-shot masked filling), and one injection
geometry (uniform difference-of-means addition), so the composition-only limit
is specific to that geometry. The usable dose range is narrow because the
evaluation degenerates at $\alpha\!\ge\!0.35$, so passes are demonstrated at
modest strength and achievable effect is capped by the evaluation rather than
necessarily by the direction. Proxies are compositional and non-mechanistic; a
pass shows the model can be pushed toward an amino-acid signature that
correlates with the property, not that mechanistic specificity is preserved.
Several validity checks were clarified after outcomes were seen, bootstrap
intervals describe stability over the fixed evaluation proteins rather than
population guarantees, and some legacy runs omit seed metadata. The separation
predictor is calibrated on nine targets from one arc and one model; whether the
same ordering holds for other models, decoders, or injection geometries is
untested.

\section*{LLM usage disclosure}
Large language models assisted with manuscript restructuring, prose editing, and
source checks, and with code and earlier drafting during the underlying study.
The authors made the scientific decisions and verified every numerical claim
against the saved artifacts.

\section*{Code and data availability}
The study uses public protein resources (UniProt, DLKcat, ProteinGym, DisProt,
IEDB, eSol) together with saved model outputs. An anonymized artifact package
with per-claim source files and seed details will be provided if the submission
system supports it. The review manuscript does not link to an identifying
repository.

\section*{Ethics declaration}
This computational study used public protein sequences and aggregated public
annotations. It involved no new human participants, clinical intervention, or
collection of personal data. Generated sequences were used only to compute
compositional proxy scores and were not synthesized or expressed.

\section*{Reproducibility statement}
Each numerical claim maps to a saved result artifact and a study record fixing
the dataset, proxy, separation, seeds, and dose grid. The intervention,
controls, and acceptance criteria were fixed before the runs they judge.

\bibliographystyle{plainnat}
\bibliography{reference}

\end{document}
